\textbf{Proof.}
Index the \(N=2n\) exposure coordinates by \(j\), put \(Z_j=D_j-\pi_j\), and let \(d_n\le\overline d n^\alpha\) be the largest closed dependency neighborhood.  If a tuple of centered coordinates has a distinct index separated from all the others, its joint moment vanishes.  Counting the remaining dependency-supported tuples and using \(|Z_j|\le1\) gives
\[
\mathbb E_\nu\left(\sum_jq_jZ_j\right)^4
\le C N^2d_n^2\lVert q\rVert_\infty^4,
\tag{1}
\]
and
\[
\left|\mathbb E_\nu
\left(\sum_jq_jZ_j\right)
\left(\sum_jr_jZ_j\right)
\left(\sum_js_jZ_j\right)\right|
\le C Nd_n^2\lVert q\rVert_\infty\lVert r\rVert_\infty\lVert s\rVert_\infty.
\tag{2}
\]
A contributing four-tuple either contains two dependency edges, giving \(O(N^2d_n^2)\) choices, or is connected, giving \(O(Nd_n^3)\le O(N^2d_n^2)\); a contributing triple has \(O(Nd_n^2)\) choices.

For the signed coordinate \(Y_j\), def:observable-coefficient-map gives
\[
\mathbb E_\nu\widehat Y_j^{\mathrm{obs}}=Y_j,
\qquad
\{\Pi^{-1}(\widehat Y^{\mathrm{obs}}-Y)\}_j
=\frac{Y_jZ_j}{\pi_j^2}.
\tag{3}
\]
Therefore
\[
G_B:=\widehat\gamma^{\mathrm{HT}}(B)-\gamma_B^*(Y)
=\sum_{j=1}^{N}u_{B,j}Z_j,
\qquad
u_{B,j}=S_B^{-1}B^{\top}\Omega e_j\frac{Y_j}{\pi_j^2}.
\tag{4}
\]
Column sparsity of \(\Omega\), the coordinate envelope, the probability-rate bound, and \(S_B\succeq n\kappa I_2\) imply
\[
\sup_{B,j}\lVert u_{B,j}\rVert_2\le K_un^{\alpha+2\beta-1}.
\tag{5}
\]
Also
\[
X_B=n^{-1}B^{\top}(D-\mathbb E_\nu D)=\sum_jx_{B,j}Z_j,
\qquad \sup_{B,j}\lVert x_{B,j}\rVert_2\le\overline a/n.
\tag{6}
\]
Applying (1) coordinatewise gives
\[
\sup_B\mathbb E_\nu\lVert G_B\rVert_2^4
\le K_Gn^{6\alpha+8\beta-2},
\qquad
\sup_B\mathbb E_\nu\lVert X_B\rVert_2^4
\le K_Xn^{2\alpha-2}.
\tag{7}
\]
Because \(\widehat\tau_B^{\mathrm{CV}}=\widehat\tau_B^{\mathrm{or}}+\Delta_B\) with \(\Delta_B=-G_B^{\top}X_B\), Cauchy--Schwarz yields
\[
\sup_B\mathbb E_\nu\Delta_B^2
\le K_\Delta n^{4\alpha+4\beta-2}.
\tag{8}
\]

Let \(T_B=\widehat\tau_B^{\mathrm{or}}-\tau\).  It is centered and has representation \(T_B=\sum_jt_{B,j}Z_j\), where
\[
t_{B,j}=n^{-1}\{Y_j/\pi_j-B_{j\cdot}\gamma_B^*(Y)\}.
\tag{9}
\]
The projector inequality
\[
\gamma_B^*(Y)^{\top}S_B\gamma_B^*(Y)
\le Y^{\top}\Pi^{-1}\Omega\Pi^{-1}Y
\]
implies
\[
\sup_B\lVert\gamma_B^*(Y)\rVert_2\le K_\gamma n^{\alpha/2+\beta},
\qquad
\sup_{B,j}|t_{B,j}|\le K_tn^{\alpha/2+\beta-1}.
\tag{10}
\]
Expanding the covariance and applying (2), (5), (6), and (10) yields
\[
\sup_B|\operatorname{Cov}_\nu(T_B,\Delta_B\mid Y)|
\le K_{\mathrm{cov}}n^{(7/2)\alpha+3\beta-2}.
\tag{11}
\]
The variance expansion with (8) and (11) proves the first uniform bound with the displayed \(R_n\).  Since \((7/2)\alpha+3\beta\le4\alpha+4\beta\), it also gives \(R_n\le K_Rn^{4\alpha+4\beta-2}\).

Completing the square for the oracle coefficient gives
\[
\operatorname{Var}_\nu(\widehat\tau_B^{\mathrm{or}}\mid Y)
=n^{-2}Y^{\top}\{\Pi^{-1}\Omega\Pi^{-1}
-\Pi^{-1}\Omega BS_B^{-1}B^{\top}\Omega\Pi^{-1}\}Y.
\tag{12}
\]
Assumption ass:iid-potential-pairs gives \(\mathbb E_{H^{\otimes n}}(YY^{\top})=Q(c_H)\).  Taking the trace expectation in (12) yields exactly \(\ell_B(c_H)\).  Integrating the first bound proves the second.  Since \(4\alpha+4\beta<1\), \(nR_n\to0\). \(\square\)